\section*{Примечания о чтениях}
\subsection*{\texttt{GELFAND-P001425}}
\par\noindent\leavevmode\phantomsection\label{GELFAND-0155-U001}\label{GELFAND-P001425:vertex-cancellation}\textbf{\texttt{GELFAND-0155-U001}} (vertex-cancellation).\nopagebreak[4]\par
Перед «Алгебра» зачёркнуто короткое начало слова; различимы петли, похожие на «Ве» и нисходящий штрих. Начало следующего слова «Вершинных» также обведено/переписано.\par
Кандидаты: Ве…; Вер…. Чтение не установлено в дипломатическом исходнике.\par
Исходный кадр 2408×3480; координаты (x0,y0,x1,y1), начало в левом верхнем углу: (285, 2500, 1060, 2655). Повторно рассмотрен увеличенный фрагмент, ×5; увеличение не добавляет исходных пикселей. Зачёркнутый фрагмент не восстановлен по тематическому контексту; разборчивое последующее слово сохранено.\par
\medskip
\subsection*{\texttt{GELFAND-P001426}}
\par\noindent\leavevmode\phantomsection\label{GELFAND-0155-U002}\label{GELFAND-P001426:frame-prefix}\textbf{\texttt{GELFAND-0155-U002}} (frame-prefix).\nopagebreak[4]\par
Непосредственно перед разборчивым «реперов» находится густо зачёркнутый короткий фрагмент. Под словом отдельно написано «(кореперов) — лучше говорить».\par
Кандидаты: ко. Чтение не установлено в дипломатическом исходнике.\par
Исходный кадр 2408×3480; координаты (x0,y0,x1,y1), начало в левом верхнем углу: (1870, 850, 2250, 945). Повторно рассмотрен увеличенный фрагмент, ×6; увеличение не добавляет исходных пикселей. Нельзя уверенно отличить отменённую приставку от переписанного начала слова. Читаемая поправка «кореперов» сохранена отдельно.\par
\medskip
\par\noindent\leavevmode\phantomsection\label{GELFAND-0155-U003}\label{GELFAND-P001426:transitivity}\textbf{\texttt{GELFAND-0155-U003}} (transitivity).\nopagebreak[4]\par
После диаграммы написано «G должно действ. на P транзитивно». Слов «в каждом слое» нет.\par
Кандидаты: на P транзитивно. Чтение не установлено в дипломатическом исходнике.\par
Исходный кадр 2408×3480; координаты (x0,y0,x1,y1), начало в левом верхнем углу: (425, 1520, 1600, 1645). Повторно рассмотрен увеличенный фрагмент, ×4; увеличение не добавляет исходных пикселей. В контексте главного расслоения транзитивность обычно относится к отдельным слоям, а не ко всему P. Возможно, это подразумеваемое сокращение записи. Уточнение не внесено в текст как авторское.\par
\medskip
\par\noindent\leavevmode\phantomsection\label{GELFAND-0155-U004}\label{GELFAND-P001426:tensor-cancellation}\textbf{\texttt{GELFAND-0155-U004}} (tensor-cancellation).\nopagebreak[4]\par
После «∃ корепер» перед запятой зачёркнут одиночный знак/очень короткое обозначение; видны петля и вертикальный штрих.\par
Кандидаты: \textbackslash{}theta; \textbackslash{}varphi. Чтение не установлено в дипломатическом исходнике.\par
Исходный кадр 2408×3480; координаты (x0,y0,x1,y1), начало в левом верхнем углу: (710, 2090, 1220, 2220). Повторно рассмотрен увеличенный фрагмент, ×5; увеличение не добавляет исходных пикселей. Кандидаты относятся только к отменённому знаку. Его значение не восстановлено по определению G-структуры.\par
\medskip
\subsection*{\texttt{GELFAND-P001427}}
\par\noindent\leavevmode\phantomsection\label{GELFAND-0155-U005}\label{GELFAND-P001427:quotient-notation}\textbf{\texttt{GELFAND-0155-U005}} (quotient-notation).\nopagebreak[4]\par
В пояснении под T\_xM написано (T\_pM/T\_p\textasciicircum{}vM); обе буквы после T имеют отчётливую форму M, хотя выше p∈P.\par
Кандидаты: T\_pM/T\_p\textasciicircum{}vM. Чтение не установлено в дипломатическом исходнике.\par
Исходный кадр 2408×3480; координаты (x0,y0,x1,y1), начало в левом верхнем углу: (905, 465, 1190, 605). Повторно рассмотрен увеличенный фрагмент, ×8; увеличение не добавляет исходных пикселей. Вероятная описка исходной записи: для ранее заданного π:P→M сопоставление с T\_xM обычно использует T\_pP/T\_p\textasciicircum{}vP. Это только математическая проверка; M в транскрипции не заменено на P.\par
\medskip
\par\noindent\leavevmode\phantomsection\label{GELFAND-0155-U006}\label{GELFAND-P001427:zero-cancellation}\textbf{\texttt{GELFAND-0155-U006}} (zero-cancellation).\nopagebreak[4]\par
Две горизонтальные черты перечёркивают строку «θ\_p на [слово] 0»; среднее слово напоминает «слое».\par
Кандидаты: слое. Чтение не установлено в дипломатическом исходнике.\par
Исходный кадр 2408×3480; координаты (x0,y0,x1,y1), начало в левом верхнем углу: (385, 490, 1160, 600). Повторно рассмотрен увеличенный фрагмент, ×6; увеличение не добавляет исходных пикселей. Смысл условия горизонтальности согласуется с кандидатом, но не доказывает чтения зачёркнутого слова.\par
\medskip
\par\noindent\leavevmode\phantomsection\label{GELFAND-0155-U007}\label{GELFAND-P001427:horizontal-cancellation}\textbf{\texttt{GELFAND-0155-U007}} (horizontal-cancellation).\nopagebreak[4]\par
После H⊂T\_pP перед тире расположен небольшой зачёркнутый фрагмент из наклонных штрихов.\par
Кандидаты: не записаны. Чтение не установлено в дипломатическом исходнике.\par
Исходный кадр 2408×3480; координаты (x0,y0,x1,y1), начало в левом верхнем углу: (1430, 2260, 1620, 2370). Повторно рассмотрен увеличенный фрагмент, ×6; увеличение не добавляет исходных пикселей. Нельзя надёжно определить исходную букву или знак; сам факт отмены передан.\par
\medskip
\par\noindent\leavevmode\phantomsection\label{GELFAND-0155-U008}\label{GELFAND-P001427:connection-cancellation}\textbf{\texttt{GELFAND-0155-U008}} (connection-cancellation).\nopagebreak[4]\par
В последней строке между «Опр.» и π′ зачёркнуто обозначение с высокой петлёй и наклонными штрихами.\par
Кандидаты: P'. Чтение не установлено в дипломатическом исходнике.\par
Исходный кадр 2408×3480; координаты (x0,y0,x1,y1), начало в левом верхнем углу: (300, 3080, 510, 3210). Повторно рассмотрен увеличенный фрагмент, ×6; увеличение не добавляет исходных пикселей. Возможное P′ не утверждается. Следующее разборчивое обозначение — π′.\par
\medskip
\subsection*{\texttt{GELFAND-P001428}}
\par\noindent\leavevmode\phantomsection\label{GELFAND-0155-U009}\label{GELFAND-P001428:decomposition-cancellation}\textbf{\texttt{GELFAND-0155-U009}} (decomposition-cancellation).\nopagebreak[4]\par
В разложении T\_pP после знака равенства перед T\_p\textasciicircum{}v зачёркнут знак с петлёй и нижними штрихами. Разборчивые слагаемые после него — T\_p\textasciicircum{}v⊕H(p′).\par
Кандидаты: не записаны. Чтение не установлено в дипломатическом исходнике.\par
Исходный кадр 2408×3480; координаты (x0,y0,x1,y1), начало в левом верхнем углу: (820, 1220, 980, 1370). Повторно рассмотрен увеличенный фрагмент, ×8; увеличение не добавляет исходных пикселей. Фрагмент рассмотрен отдельно. Ни исходную букву, ни возможный индекс надёжно восстановить нельзя. Вертикальная стрелка i\_p\textasciicircum{}\{-1\} и нижняя g не относятся к отмене и переданы в схеме.\par
\medskip
\subsection*{\texttt{GELFAND-P001429}}
\par\noindent\leavevmode\phantomsection\label{GELFAND-0155-U010}\label{GELFAND-P001429:speaker-cancellation}\textbf{\texttt{GELFAND-0155-U010}} (speaker-cancellation).\nopagebreak[4]\par
Между вопросом о гильбертовом пространстве и подчёркнутой фамилией находится короткая зачёркнутая строка, начинающаяся знаком, похожим на «З». Видны следующие соединённые штрихи и нижняя петля.\par
Кандидаты: Зам…. Чтение не установлено в дипломатическом исходнике.\par
Исходный кадр 2408×3480; координаты (x0,y0,x1,y1), начало в левом верхнем углу: (305, 1000, 555, 1090). Повторно рассмотрен увеличенный фрагмент, ×6; увеличение не добавляет исходных пикселей. Кандидат не дополнен до фамилии по следующему заголовку.\par
\medskip
\par\noindent\leavevmode\phantomsection\label{GELFAND-0155-U011}\label{GELFAND-P001429:equation-number}\textbf{\texttt{GELFAND-0155-U011}} (equation-number).\nopagebreak[4]\par
У определения Z конечный вид номера читается как (4); под правой частью сохраняются более ранние округлые штрихи, похожие на 3.\par
Кандидаты: 3. Чтение не установлено в дипломатическом исходнике.\par
Исходный кадр 2408×3480; координаты (x0,y0,x1,y1), начало в левом верхнем углу: (150, 2440, 265, 2575). Повторно рассмотрен увеличенный фрагмент, ×8; увеличение не добавляет исходных пикселей. Конечная нумерация (4) сохранена; первоначальный номер остаётся кандидатом.\par
\medskip
\par\noindent\leavevmode\phantomsection\label{GELFAND-0155-U012}\label{GELFAND-P001429:algebra-cancellation}\textbf{\texttt{GELFAND-0155-U012}} (algebra-cancellation).\nopagebreak[4]\par
После «ассоц. алгебру;» перед «с образующи…» зачёркнут короткий фрагмент. Выносная строка над ним читается «(с континуумом)».\par
Кандидаты: не записаны. Чтение не установлено в дипломатическом исходнике.\par
Исходный кадр 2408×3480; координаты (x0,y0,x1,y1), начало в левом верхнем углу: (1125, 2815, 1970, 2990). Повторно рассмотрен увеличенный фрагмент, ×5; увеличение не добавляет исходных пикселей. Первоначальный короткий фрагмент не восстановлен. Разборчивая надстрочная вставка сохранена отдельно.\par
\medskip
\par\noindent\leavevmode\phantomsection\label{GELFAND-0155-U013}\label{GELFAND-P001429:algebra-ending}\textbf{\texttt{GELFAND-0155-U013}} (algebra-ending).\nopagebreak[4]\par
Слово в строке после «с» имеет основу «образующи…»; последние штрихи переписаны, видны элементы окончания «ми» и наложенного «х». Далее A\_i(x).\par
Кандидаты: образующими; образующих. Чтение не установлено в дипломатическом исходнике.\par
Исходный кадр 2408×3480; координаты (x0,y0,x1,y1), начало в левом верхнем углу: (1125, 2815, 1970, 2990). Повторно рассмотрен увеличенный фрагмент, ×5; увеличение не добавляет исходных пикселей. Вставка «с континуумом» делает возможным грамматическое изменение окончания, но окончательный вариант не принят автоматически. В теле приведены обе формы.\par
\medskip
\subsection*{\texttt{GELFAND-P001430}}
\par\noindent\leavevmode\phantomsection\label{GELFAND-0155-U014}\label{GELFAND-P001430:integral-cancellation}\textbf{\texttt{GELFAND-0155-U014}} (integral-cancellation).\nopagebreak[4]\par
Между «Рассм.» и «интеграл» зачёркнуто слово с двумя нижними петлями; читаемые остатки напоминают «оператор».\par
Кандидаты: оператор. Чтение не установлено в дипломатическом исходнике.\par
Исходный кадр 2408×3480; координаты (x0,y0,x1,y1), начало в левом верхнем углу: (400, 1115, 610, 1220). Повторно рассмотрен увеличенный фрагмент, ×6; увеличение не добавляет исходных пикселей. Кандидат согласуется с последующей строкой о линейном операторе, однако зачёркивание закрывает часть букв.\par
\medskip
\par\noindent\leavevmode\phantomsection\label{GELFAND-0155-U015}\label{GELFAND-P001430:epsilon-first-index}\textbf{\texttt{GELFAND-0155-U015}} (epsilon-first-index).\nopagebreak[4]\par
В (7) второй нижний индекс у ε густо переписан. Конечные штрихи похожи на λ; под ними заметны дополнительные штрихи.\par
Кандидаты: \textbackslash{}nu. Чтение не установлено в дипломатическом исходнике.\par
Исходный кадр 2408×3480; координаты (x0,y0,x1,y1), начало в левом верхнем углу: (690, 1240, 970, 1465). Повторно рассмотрен увеличенный фрагмент, ×6; увеличение не добавляет исходных пикселей. В тексте оставлен конечный λ, подтверждаемый разборчивой выноской ε\_\{μλ\}=−ε\_\{λμ\}. Кандидат ν относится только к первоначальному, перекрытому знаку.\par
\medskip
\par\noindent\leavevmode\phantomsection\label{GELFAND-0155-U016}\label{GELFAND-P001430:epsilon-repeat}\textbf{\texttt{GELFAND-0155-U016}} (epsilon-repeat).\nopagebreak[4]\par
В повторном интегральном равенстве после определения тензора энергии-импульса написано ε\_\{μν\}, затем dx\textasciicircum{}λ, справа ∂/∂y\textasciicircum{}ν. Второй индекс у ε имеет форму ν, а не λ.\par
Кандидаты: \textbackslash{}varepsilon\_\{\textbackslash{}mu\textbackslash{}nu\}. Чтение не установлено в дипломатическом исходнике.\par
Исходный кадр 2408×3480; координаты (x0,y0,x1,y1), начало в левом верхнем углу: (430, 1910, 1330, 2070). Повторно рассмотрен увеличенный фрагмент, ×6; увеличение не добавляет исходных пикселей. Вероятная исходная индексная описка: сопоставление с (7) и свёрткой по dx\textasciicircum{}λ указывает на ожидаемый ε\_\{μλ\}, но такое исправление не внесено. Нижний ν сохранён; математическое согласование остаётся задачей рецензента.\par
\medskip
\subsection*{\texttt{GELFAND-P001431}}
\par\noindent\leavevmode\phantomsection\label{GELFAND-0155-U017}\label{GELFAND-P001431:tensor-components}\textbf{\texttt{GELFAND-0155-U017}} (tensor-components).\nopagebreak[4]\par
В (10) zz и z̄z̄ стоят именно сверху: T=T\textasciicircum{}\{zz\}, T̄=T\textasciicircum{}\{z̄z̄\}; Θ=T\textasciicircum{}\{zz̄\}=0. Далее (11) утверждает ∂\_\{z̄\}T=0 и ∂\_z T̄=0.\par
Кандидаты: T=T\textasciicircum{}\{zz\}, \textbackslash{}bar T=T\textasciicircum{}\{\textbackslash{}bar z\textbackslash{}bar z\}. Чтение не установлено в дипломатическом исходнике.\par
Исходный кадр 2408×3480; координаты (x0,y0,x1,y1), начало в левом верхнем углу: (335, 355, 945, 610). Повторно рассмотрен увеличенный фрагмент, ×6; увеличение не добавляет исходных пикселей. Верхнее положение индексов подтверждено увеличением. Если понимать их как обычные контравариантные компоненты и буквально использовать (6) предыдущего кадра, при нулевой смешанной компоненте получается противоположная голоморфная зависимость. Возможно, автор использует неоговорённое соглашение о компонентах либо допустил описку. Нижние индексы и перестановка T/T̄ не подставлены. Для проверки сопоставить также (11), область (330,1070,1840,1470).\par
\medskip
\subsection*{\texttt{GELFAND-P001432}}
\par\noindent\leavevmode\phantomsection\label{GELFAND-0155-U018}\label{GELFAND-P001432:laurent-factor}\textbf{\texttt{GELFAND-0155-U018}} (laurent-factor).\nopagebreak[4]\par
В (16) второй множитель записан (z-z\_2)\textasciicircum{}\{-3\}, без индекса 1 при первом z; в первом множителе написано (z\_1-z\_2)\textasciicircum{}\{-4\}. Индексы T\_\{-4\} и T\_\{-3\} содержат минусы.\par
Кандидаты: (z-z\_2)\textasciicircum{}\{-3\}. Чтение не установлено в дипломатическом исходнике.\par
Исходный кадр 2408×3480; координаты (x0,y0,x1,y1), начало в левом верхнем углу: (260, 390, 1400, 540). Повторно рассмотрен увеличенный фрагмент, ×6; увеличение не добавляет исходных пикселей. Предыдущее разложение (15) использует переменные z\_1,z\_2, так что отсутствие индекса может быть опиской; индекс не дописан. Первый коэффициент T\_\{-4\} проверен отдельно на фрагменте (630,380,900,535), ×8: минус различим.\par
\medskip
\par\noindent\leavevmode\phantomsection\label{GELFAND-0155-U019}\label{GELFAND-P001432:lower-cancellation}\textbf{\texttt{GELFAND-0155-U019}} (lower-cancellation).\nopagebreak[4]\par
Между «зависит» и «ни от какой теории» зачёркнут короткий фрагмент; сохранившиеся штрихи не дают уверенного чтения.\par
Кандидаты: от; не. Чтение не установлено в дипломатическом исходнике.\par
Исходный кадр 2408×3480; координаты (x0,y0,x1,y1), начало в левом верхнем углу: (475, 3180, 620, 3275). Повторно рассмотрен увеличенный фрагмент, ×8; увеличение не добавляет исходных пикселей. Оба варианта остаются кандидатами; зачёркнутый фрагмент не опущен.\par
\medskip
\subsection*{\texttt{GELFAND-P001433}}
\par\noindent\leavevmode\phantomsection\label{GELFAND-0155-U020}\label{GELFAND-P001433:conformal-factor}\textbf{\texttt{GELFAND-0155-U020}} (conformal-factor).\nopagebreak[4]\par
Перед (x) в правой части уравнения ∂\_με\_ν+∂\_νε\_μ=δ\_\{μν\}…(x) написан петлеобразный знак с загнутым влево нижним штрихом.\par
Кандидаты: \textbackslash{}rho(x); g(x). Чтение не установлено в дипломатическом исходнике.\par
Исходный кадр 2408×3480; координаты (x0,y0,x1,y1), начало в левом верхнем углу: (980, 3150, 1240, 3290). Повторно рассмотрен увеличенный фрагмент, ×8; увеличение не добавляет исходных пикселей. Форма согласуется с рукописным ρ, но латинское g не исключено. Уравнение конформного преобразования само по себе не определяет выбор буквы. В строке принят кандидат ρ, а не утверждено окончательное чтение.\par
\medskip
\par\noindent\leavevmode\phantomsection\label{GELFAND-0155-U021}\label{GELFAND-P001433:trace-equality}\textbf{\texttt{GELFAND-0155-U021}} (trace-equality).\nopagebreak[4]\par
Последняя строка имеет вид T\textasciicircum{}\{z z̄\}=T\textasciicircum{}μ\_μ — след; =0. Дополнительного числового коэффициента в равенстве нет.\par
Кандидаты: T\textasciicircum{}\{z\textbackslash{}bar z\}=T\textasciicircum{}\textbackslash{}mu\_\textbackslash{}mu. Чтение не установлено в дипломатическом исходнике.\par
Исходный кадр 2408×3480; координаты (x0,y0,x1,y1), начало в левом верхнем углу: (430, 3270, 1125, 3405). Повторно рассмотрен увеличенный фрагмент, ×6; увеличение не добавляет исходных пикселей. Смешанная компонента прямо отождествлена со следом. Числовой коэффициент в таком отождествлении зависит от соглашения о комплексных координатах и тензорных компонентах; при обычном преобразовании компонентов с объявленным ds²=dz dz̄ это требует проверки. В частном случае нулевого следа оба выражения обращаются в нуль. Источник воспроизведён без внесения коэффициента.\par
Новая отдельная проверка: при обычном контравариантном преобразовании симметричного тензора и объявленных комплексных координатах равенство смешанной компоненты вещественному евклидову следу верно без дополнительного коэффициента. Историческое замечание сохранено; вычисление дано отдельно.\par
\medskip
\subsection*{\texttt{GELFAND-P001434}}
\par\noindent\leavevmode\phantomsection\label{GELFAND-0155-U022}\label{GELFAND-P001434:descendant-bracket}\textbf{\texttt{GELFAND-0155-U022}} (descendant-bracket).\nopagebreak[4]\par
Строка под (26) записана как L\_\{n\_N\}…L\_\{n\_1\}φ\_α=[φ\_α], со знаком равенства; далее стоит n\_i<0. Фигурных скобок, знака линейной оболочки или произведения операторов L̄ здесь нет.\par
Кандидаты: L\_\{n\_N\}\textbackslash{}ldots L\_\{n\_1\}\textbackslash{}phi\_\textbackslash{}alpha=[\textbackslash{}phi\_\textbackslash{}alpha]. Чтение не установлено в дипломатическом исходнике.\par
Исходный кадр 2408×3480; координаты (x0,y0,x1,y1), начало в левом верхнем углу: (255, 590, 835, 715). Повторно рассмотрен увеличенный фрагмент, ×8; увеличение не добавляет исходных пикселей. Вероятно, это сокращённое обозначение семейства потомков, но явного определения скобок источник не даёт. Равенство не заменено на принадлежность или линейную оболочку. Индекс у правого L после увеличения читается n\_1. В следующей строке прямой суммы Φ\_α имеет верхнюю и нижнюю засечки, в отличие от φ\_α здесь; прописное Φ сохранено, фрагмент (540,915,755,1020), ×8.\par
\medskip
\par\noindent\leavevmode\phantomsection\label{GELFAND-0155-U023}\label{GELFAND-P001434:barred-coordinate}\textbf{\texttt{GELFAND-0155-U023}} (barred-coordinate).\nopagebreak[4]\par
В нижней паре замен справа во второй строке написано λ̄z: над λ черта есть, над следующим z её не видно. Слева второй строки написано z̄.\par
Кандидаты: \textbackslash{}bar z\textbackslash{}to\textbackslash{}bar\textbackslash{}lambda z. Чтение не установлено в дипломатическом исходнике.\par
Исходный кадр 2408×3480; координаты (x0,y0,x1,y1), начало в левом верхнем углу: (215, 1270, 500, 1460). Повторно рассмотрен увеличенный фрагмент, ×8; увеличение не добавляет исходных пикселей. Возможно, пропущена черта над правым z. Не исправлено на λ̄z̄. Степенной множитель справа проверен отдельно: λ\textasciicircum{}\{Δ\_α\} λ̄\textasciicircum{}\{Δ̄\_α\}; короткий горизонтальный штрих над второй λ — черта, а не минус показателя. Области (620,1260,1240,1430), ×8, и (900,1270,1080,1410), ×10.\par
\medskip
